\documentclass[11pt,a4paper]{article}
\usepackage[T1]{fontenc}
\usepackage[utf8]{inputenc}
\usepackage[ngerman]{babel}
\usepackage{lmodern}
\usepackage[margin=2.2cm]{geometry}
\usepackage{microtype}
\usepackage{enumitem}
\usepackage[hidelinks]{hyperref}

\setlength{\parindent}{0pt}
\setlength{\parskip}{0.4em}
\setlist[itemize]{leftmargin=*,itemsep=0.15em,topsep=0.2em}
\setlist[enumerate]{leftmargin=*,itemsep=0.35em,topsep=0.2em}
\emergencystretch=3em

\title{Vorbereitung des Betreuungsgesprächs\\
\large Literaturrecherche und Versuchsplanung}
\author{Bachelorarbeit: Tabular Foundation Models bei fehlenden Merkmalswerten}
\date{Termin: 27. November 2026 \quad|\quad Arbeitsstand: 25. November 2026}

\begin{document}
\maketitle

\textbf{Zweck:} Zwischenstand seit Abgabe des Exposés vorstellen und
Rückmeldung zum weiteren Vorgehen einholen. Die Recherche und die folgenden
Planungen sind noch nicht abgeschlossen.

\section*{1. Bisheriger Stand}

\textbf{Literaturrecherche (FF1).}
\begin{itemize}
  \item Erste Suchläufe in Google Scholar, arXiv und IEEE Xplore sind
  durchgeführt und im Suchprotokoll dokumentiert. Scopus war trotz VPN nur
  als Preview zugänglich; laut DBIS besteht derzeit kein Uni-Zugriff.
  \item Suchanfragen, Trefferzahlen und -- soweit bekannt -- Sichtungsbereiche
  werden je Suchlauf festgehalten. Bei Google Scholar wurden in S1a--S1c,
  S2a, S3a und S4a jeweils die ersten 50 Treffer gesichtet; der Sichtungsbereich
  der explorativen Suche S3b ist nicht dokumentiert.
  \item Treffer werden in Zotero gesammelt. Ein erster DOI-/Titelabgleich der
  lokalen RIS-Exporte ist dokumentiert; die Prüfung in Zotero sowie die
  systematische Auswahl und Volltextprüfung stehen noch aus. Die Konzeptmatrix
  ist vorbefüllt, aber noch am Volltext zu prüfen.
  \item Zaimoğlu, Bi und Jang sind Beispiele für vorläufig interessante
  Kandidaten. Ihre endgültige   Einordnung erfordert Volltextprüfung. Die erste Datenbanksuchrunde ist
  vorläufig beendet; weitere Ergänzungen sollen mit dem Betreuer geklärt
  werden. Vorwärts-/Rückwärtssuche steht noch aus.
\end{itemize}

\textbf{Technischer Stand (FF2).}
\begin{itemize}
  \item TabPFN-3.5 und LightGBM laufen im technischen Pilot auf WDBC; der
  Pilot wurde auf der PC-GPU wiederholt. Ein zusätzlicher Laufzeittest auf
  Spambase liegt vor.
  \item Beide Tests dienen der Machbarkeit, nicht als Ergebnis des
  Hauptversuchs. Der Spambase-Test umfasst einen Split und nur MCAR; die
  vollständige Hauptversuchspipeline sowie MAR und MNAR sind noch nicht
  implementiert beziehungsweise getestet.
  \item Vorgesehen sind MCAR, MAR und MNAR; die konkreten Regeln und der
  Umfang je Mechanismus sind offen. Interne Verarbeitung und Median-Imputation
  bilden den Kern; Median mit Missing-Indikatoren ist als Zusatzanalyse
  vorgesehen.
\end{itemize}

\textbf{Schreiben und Zeit:} Das LaTeX-Kapitelgerüst steht, Kapiteltext ist
noch nicht geschrieben. Abgabe laut Projektstand: 16. Februar 2027. Das
Seitenbudget von ungefähr 40 Seiten ist bisher keine bestätigte Vorgabe.

\section*{2. Suchstrategie und Suchbegriffe}

Die Suche verbindet Modellansätze und fehlende Merkmalswerte. Einzelne
Suchläufe zu nur einem Modellansatz bleiben wichtig; direkte Vergleiche sind
nicht die einzige Einschlussmöglichkeit. Grundlagen und Methodikquellen
werden bei Bedarf getrennt ergänzt.

\begin{itemize}
  \item \textbf{Modelle:} TabPFN, tabular foundation models, gradient boosting,
  LightGBM und XGBoost.
  \item \textbf{Fehlende Werte und Behandlung:} missing values, missing data,
  missingness, imputation; ergänzend MCAR, MAR und MNAR.
  \item \textbf{Aufgabe und Vergleich:} classification, prediction und
  benchmark.
\end{itemize}

Die folgenden Angaben zeigen die bisher verwendeten Suchideen.
Exakte Eingaben, Felder, Filter und noch ungeklärte Angaben stehen im
Suchprotokoll; eine Vorlage ist kein Nachweis einer durchgeführten Suche.

\subsection*{Google Scholar: durchgeführte Varianten}
\begin{itemize}
  \item S1a: \verb|"TabPFN" "missing values"|
  \item S1b: \verb|"tabular foundation models" "missing data"|
  \item S1c: \verb|"TabPFN" imputation|
  \item S2a: \verb|"TabPFN" "gradient boosting" "missing values"|
  \item S3a: \verb|"gradient boosting" "missing values"|
  \item S4a: \verb|"missing data" imputation classification benchmark|
\end{itemize}
S3b wurde zusätzlich explorativ versucht; fehlende Suchdetails sind im
Suchprotokoll ausdrücklich offengelassen.

\subsection*{arXiv: durchgeführte Varianten}
\begin{itemize}
  \item S1a: \verb|TabPFN| AND \verb|"missing values"|
  \item S1b: \verb|"tabular foundation models"| AND \verb|"missing data"|
  \item S1c: \verb|TabPFN| AND \verb|imputation|
  \item S2a: \verb|"TabPFN"|, \verb|"gradient boosting"| und
  \verb|"missing values"|; die zweite Verknüpfung ist nicht dokumentiert.
  \item S3a: \verb|"gradient boosting"| AND \verb|"missing values"|
  \item S4a: \verb|"missing data"| AND \verb|imputation| AND
  \verb|classification| AND \verb|benchmark|
  \item S5a: \verb|"missing data"| AND \verb|classification| AND
  \verb|benchmark|
\end{itemize}
Die Suchfelder waren laut Nutzer auf „All fields“ eingestellt; gemeinsame
Einstellungen und Abweichungen sind im Suchprotokoll festgehalten.

\subsection*{IEEE Xplore: durchgeführte und geplante Varianten}
\begin{itemize}
  \item S1a: \verb|"TabPFN"| AND \verb|"missing values"|
  \item S1b: \verb|"tabular foundation models"| AND \verb|"missing data"|
  (beide Felder „All Metadata“; keine Treffer).
  \item S1c: \verb|"TabPFN"| AND \verb|"imputation"|
  \item S3a: \verb|"gradient boosting"| AND \verb|"missing values"|
  \item Für S2a und S4a liegen Exporte vor; die genauen Eingaben sind noch
  nicht bestätigt und im Suchprotokoll entsprechend gekennzeichnet.
  \item S1b wurde zusätzlich mit Alternativen für TabPFN beziehungsweise
  Tabular Foundation Models und für \verb|"missing data"| beziehungsweise
  \verb|"missing values"| erweitert.
  \item E1 ergänzte diesen Ansatz um \verb|"missingness"| und
  \verb|"imputation"|: 10 Treffer laut Nutzer. E2 verknüpfte
  \verb|"gradient boosting"|, \verb|"LightGBM"| oder \verb|"XGBoost"| mit
  \verb|"missing values"|, \verb|"missing data"| oder \verb|"missingness"|:
  2\,525 Treffer laut Nutzer.
\end{itemize}

\textbf{Weitere Ergänzungen} sind bis zur Rücksprache vorläufig
zurückgestellt. Bei Durchführung werden sie als neue Suchläufe dokumentiert,
nicht rückwirkend in alte Läufe eingetragen. Weder eine Nulltreffersuche noch
die Sichtung der ersten 50 Ergebnisse beweist Vollständigkeit.

\section*{3. Fragen für das Gespräch}

\textbf{Priorität: Literaturrecherche}
\begin{enumerate}
  \item Sind Google Scholar, arXiv und IEEE Xplore als Ausgangspunkt
  angemessen? Ist eine weitere Datenbank nötig, wenn Scopus nicht zugänglich
  ist?
  \item Sind die Suchblöcke und bisherigen OR-Ergänzungen ausreichend breit?
  Wie soll bei vielen Treffern eine Begrenzung auf die ersten 50 begründet
  werden, und wann ist eine Zitationssuche ausreichend?
  \item Sind 5--10 Kernpaper insgesamt eine sinnvolle Orientierung? Wie
  umfangreich soll die Literaturanalyse im Verhältnis zum Experiment sein?
\end{enumerate}

\textbf{Priorität: Experiment}
\begin{enumerate}[resume]
  \item Welcher Umfang und welche Regeln für MAR und MNAR sind neben der
  MCAR-Referenz für eine Bachelorarbeit angemessen?
  \item Sind die vier numerischen Datensatzkandidaten als Ausgangspunkt
  ausreichend? Wie viele Splits und Maskierungswiederholungen sind sinnvoll?
  \item Reicht die geplante deskriptive Auswertung der gepaarten
  Leistungsabstände, oder werden statistische Tests erwartet?
  \item Passt der geplante Vergleich von interner Verarbeitung und
  Median-Imputation mit Missing-Indikatoren als Zusatzanalyse?
\end{enumerate}

\textbf{Form und Organisation}
\begin{enumerate}[resume]
  \item Welcher Seitenumfang ist verbindlich, und was zählt dazu? Welches
  Verhältnis von Literatur- und Praxisteil wird erwartet?
  \item Ist der IEEE-Zitierstil korrekt, obwohl die Vorlage LNCS verwendet?
  Wie soll die KI-Nutzung im Anhang angegeben werden?
  \item Ist der vorläufige Zeitplan bis zur Abgabe realistisch?
\end{enumerate}

\textbf{Arbeitsgrundlagen:} \texttt{literature/notes/Suchprotokoll.md}
(Suchläufe und Trefferlisten), \texttt{literature/notes/Konzeptmatrix.csv},
\texttt{experiments/Pilotprotokoll.md}, \texttt{PROJEKTSTAND.md} und Exposé.
Aktualisierungen nach weiteren Suchläufen hier ergänzen; das Suchprotokoll
bleibt die detaillierte Dokumentation.

\end{document}
